\documentclass[12pt,a4paper]{article}
\usepackage[a4paper,margin=2.4cm]{geometry}
\usepackage[T2A]{fontenc}
\usepackage[utf8]{inputenc}
\usepackage[russian]{babel}
\usepackage{amsmath,amssymb,mathtools}
\usepackage{array,booktabs,longtable}
\usepackage{xcolor}
\usepackage{tikz}
\usepackage{hyperref}
\hypersetup{colorlinks=true,linkcolor=blue!60!black,urlcolor=blue!60!black}
\setlength{\parindent}{0pt}
\setlength{\parskip}{0.55em}
\renewcommand{\arraystretch}{1.45}
\linespread{1.08}
\sloppy

\begin{document}
\thispagestyle{empty}
\vspace*{\fill}
\begin{center}
{\LARGE\bfseries Ревью домашней работы\par}
\vspace{2.0cm}
{\large Автор: Лёвин Артём Александрович\par}
\vspace{0.8cm}
{\large 17 октября 2026 года\par}
\end{center}
\vspace*{\fill}
\begin{center}
\begin{tikzpicture}[x=1cm,y=1cm]
\draw[blue!35!black,line width=0.55pt]
 (0,0) .. controls (2.0,0.6) and (3.9,-0.6) .. (6.1,0)
 .. controls (8.0,0.5) and (9.5,-0.3) .. (11.1,0);
\end{tikzpicture}
\end{center}
\newpage

\section*{Сводная таблица}
\small
\setlength{\tabcolsep}{4pt}
\begin{longtable}{@{}>{\raggedright\arraybackslash}p{0.105\linewidth}>{\raggedright\arraybackslash}p{0.19\linewidth}>{\raggedright\arraybackslash}p{0.295\linewidth}>{\raggedright\arraybackslash}p{0.14\linewidth}>{\raggedright\arraybackslash}p{0.17\linewidth}@{}}
\toprule
\textbf{№ задания} & \textbf{Тема} & \textbf{Суть ошибки} & \textbf{Ваш ответ} & \textbf{Правильный ответ}\\
\midrule
\endfirsthead
\toprule
\textbf{№ задания} & \textbf{Тема} & \textbf{Суть ошибки} & \textbf{Ваш ответ} & \textbf{Правильный ответ}\\
\midrule
\endhead
\hyperref[task:1]{1} & Действия со степенями & Отрицательный показатель учтён с неверным знаком; далее неверно вычислена разность показателей. & $256$ & $8$\\
\midrule
\hyperref[task:5]{5} & Степенная зависимость давления от объёма & После подстановки отношений давления и объёма перепутаны основание и значение степени. & $\frac{1}{2}$ & $2$\\
\bottomrule
\end{longtable}
\normalsize

\newpage
\thispagestyle{empty}
\vspace*{\fill}
\begin{center}
{\LARGE\bfseries Разбор ошибок}
\end{center}
\vspace*{\fill}
\newpage

\phantomsection\label{task:1}
\section*{Задание № 1}
\textbf{Текст задания.} Вычислите:
\[
\frac{2^8\cdot 2^{-3}}{2^2}.
\]

\textbf{Ваше решение.} В записи выполнен переход:
\[
\frac{2^8\cdot 2^{-3}}{2^2}
=2^{11-2}=2^8=256.
\]
\textbf{Ваш ответ:} $256$.

\textcolor{red}{\textbf{Ошибка:}} при умножении степеней с одинаковым основанием неверно учтён отрицательный показатель. Дополнительно в следующем переходе неправильно вычислена разность показателей.

\textbf{Где именно возникает ошибка.}

\emph{Последний правильный шаг} --- исходная дробь
\[
\frac{2^8\cdot 2^{-3}}{2^2}.
\]
\emph{Первый неправильный шаг} --- замена этой дроби выражением $2^{11-2}$. Число $11$ не может получиться при сложении показателей $8$ и $-3$:
\[
8+(-3)=8-3=5,\qquad \text{а не }11.
\]
Кроме того, даже для ошибочно записанного показателя разность $11-2$ равна $9$, а не $8$. Это отдельная арифметическая неточность. Равенство $2^8=256$ само по себе верно, но не относится к значению исходного выражения.

\textbf{Почему это неверно.} Если основания одинаковы, при умножении степеней \emph{складываются показатели}. Отрицательный показатель сохраняет знак. При делении степеней с одинаковым ненулевым основанием \emph{показатель знаменателя вычитается} из показателя числителя. Нельзя заменить $-3$ на $+3$: это меняет множитель $2^{-3}=\frac18$ на совершенно другое число $2^3=8$.

\textcolor{blue!60!black}{\textbf{Ключевая идея:}} для $c\ne0$ действуют правила
\[
c^m\cdot c^n=c^{m+n},\qquad
\frac{c^m}{c^n}=c^{m-n},\qquad
c^{-n}=\frac{1}{c^n}.
\]
Сначала применяем правило для произведения, затем для частного; на каждом шаге сохраняем знаки показателей.

\textbf{Исправление от последнего правильного шага.}
\[
\frac{2^8\cdot 2^{-3}}{2^2}
=2^{8+(-3)-2}
=2^{8-3-2}
=2^3=8.
\]

\textbf{Полное правильное решение.}

\emph{Шаг 1.} Перемножим степени двойки в числителе. Показатели складываются:
\[
2^8\cdot 2^{-3}=2^{8+(-3)}=2^5.
\]
\emph{Шаг 2.} Разделим полученную степень на $2^2$. Теперь вычитаем показатель знаменателя:
\[
\frac{2^5}{2^2}=2^{5-2}=2^3.
\]
\emph{Шаг 3.} Вычислим степень:
\[
2^3=2\cdot2\cdot2=8.
\]

\textbf{Проверка.} Посчитаем исходную дробь через обычные числовые значения:
\[
2^8=256,\qquad 2^{-3}=\frac18,\qquad 2^2=4;
\]
\[
\frac{256\cdot\frac18}{4}=\frac{32}{4}=8.
\]
Оба способа приводят к одному результату.

\textcolor{teal!60!black}{\textbf{Итог:}} правильный ответ --- $\boxed{8}$.

\textbf{Задача на закрепление.} Вычислите:
\[
\frac{3^7\cdot3^{-4}}{3^2}.
\]

\newpage
\phantomsection\label{task:5}
\section*{Задание № 5}
\textbf{Текст задания.} В процессе $pV^a=\mathrm{const}$ объём увеличился в $4$ раза, а давление уменьшилось в $16$ раз. Найдите $a$.

\textbf{Ваше решение.} Верно записана связь двух состояний и соотношения величин:
\[
p_1V_1^a=p_2V_2^a,\qquad
p_1=16p_2,\qquad V_2=4V_1.
\]
После подстановки и сокращения получено равенство
\[
16^a=4,
\]
из которого найдено $a=\frac12$.

\textbf{Ваш ответ:} $\frac12$.

\textcolor{red}{\textbf{Ошибка:}} при переходе от уравнения для двух состояний к уравнению с одной степенью перепутаны основание и значение степени. Правильное равенство имеет вид $4^a=16$, а не $16^a=4$.

\textbf{Где именно возникает ошибка.}

\emph{Последний правильный шаг} --- подстановка данных в закон процесса:
\[
16p_2V_1^a=p_2(4V_1)^a.
\]
\emph{Первый неправильный шаг} --- получение уравнения $16^a=4$ после сокращения. Коэффициент $16$ относится к давлению и не возводится в степень $a$; во вторую степень входит множитель $4$ при объёме.

\textbf{Почему это неверно.} По свойству степени произведения
\[
(4V_1)^a=4^a V_1^a.
\]
Значит, справа от равенства возникает множитель $4^a$, а слева остаётся множитель $16$. После деления обеих частей уравнения на общий ненулевой множитель $p_2V_1^a$ получаем $16=4^a$. Перестановка чисел местами без математического основания меняет уравнение и его решение. Давление и объём в данной физической задаче положительны, поэтому указанное сокращение допустимо.

\textcolor{blue!60!black}{\textbf{Ключевая идея:}} изменение каждой физической величины нужно подставлять \emph{в то место формулы, где эта величина стоит}. Если объём увеличился в $k$ раз, то $V_2=kV_1$ и $V_2^a=k^a V_1^a$; если давление уменьшилось в $m$ раз, то $p_2=\frac{p_1}{m}$.

\textbf{Исправление от последнего правильного шага.}
\begin{align*}
16p_2V_1^a&=p_2(4V_1)^a,\\
16p_2V_1^a&=4^a p_2V_1^a,\\
16&=4^a,\\
4^2&=4^a,\\
a&=2.
\end{align*}
В предпоследнем переходе использовано равенство $16=4^2$. При основании $4>1$ одинаковым значениям степени соответствуют одинаковые показатели.

\textbf{Полное правильное решение.}

\emph{Шаг 1.} Обозначим начальные давление и объём через $p_1$ и $V_1$, конечные --- через $p_2$ и $V_2$. По условию процесса произведение $pV^a$ сохраняется:
\[
p_1V_1^a=p_2V_2^a.
\]
\emph{Шаг 2.} Переведём слова условия в равенства. Объём стал вчетверо больше, а давление --- в шестнадцать раз меньше:
\[
V_2=4V_1,\qquad p_2=\frac{p_1}{16}.
\]
\emph{Шаг 3.} Подставим оба выражения в закон процесса:
\[
p_1V_1^a=\frac{p_1}{16}(4V_1)^a.
\]
\emph{Шаг 4.} Раскроем степень произведения, не теряя показатель $a$:
\[
p_1V_1^a=\frac{p_1}{16}\cdot4^aV_1^a.
\]
\emph{Шаг 5.} Разделим на положительное число $p_1V_1^a$:
\[
1=\frac{4^a}{16},\qquad 4^a=16.
\]
\emph{Шаг 6.} Представим правую часть как степень четвёрки и приравняем показатели:
\[
4^a=4^2,\qquad a=2.
\]

\textbf{Проверка.} Если $a=2$, то после изменения величин
\[
p_2V_2^2
=\frac{p_1}{16}(4V_1)^2
=\frac{p_1}{16}\cdot16V_1^2
=p_1V_1^2.
\]
Произведение сохранилось --- это совпадает с условием. Для найденного в работе значения $a=\frac12$ произведение, напротив, изменилось бы:
\[
p_2V_2^{1/2}
=\frac{p_1}{16}(4V_1)^{1/2}
=\frac18\,p_1V_1^{1/2},
\]
то есть постоянным не осталось бы.

\textcolor{teal!60!black}{\textbf{Итог:}} правильный ответ --- $\boxed{a=2}$.

\textbf{Задача на закрепление.} В процессе $pV^a=\mathrm{const}$ объём увеличился в $3$ раза, а давление уменьшилось в $81$ раз. Найдите $a$.

\vfill
\begin{center}
\footnotesize Лёвин Артём Александрович эксклюзивно для Марка Гукина
\end{center}
\end{document}
